\section{The Conference Service}
\label{sec:ch02-the-conference-service}

The example the rest of this book uses is a conference: sessions, the people
who give them, and later the ratings attendees leave. It is small enough to
print and it has the one property this book needs: everybody wants sessions
sorted by rating, and the service that owns sessions holds no ratings.
Chapter~\ref{ch:cross-seam-paging} is where that bites. For now everything is
in one process and in memory, and chapter~\ref{ch:data-without-n-plus-1} puts
it in SQLite.

Two records carry the domain. \code{Session.cs} first:

\begin{minted}{csharp}
namespace Sessions;

/// <summary>
/// A talk on the conference program.
/// </summary>
/// <param name="Id">Stable identifier for this session.</param>
/// <param name="Title">The line that goes in the program.</param>
/// <param name="Abstract">The longer text, if there is one.</param>
/// <param name="SpeakerId">Who gives it. Not part of the schema.</param>
/// <param name="StartsAt">When it starts, in conference time.</param>
/// <param name="DurationMinutes">How long the slot is.</param>
public sealed record Session(
    int Id,
    string Title,
    string? Abstract,
    [property: GraphQLIgnore] int SpeakerId,
    DateTimeOffset StartsAt,
    int DurationMinutes);
\end{minted}

Then \code{Speaker.cs}:

\begin{minted}{csharp}
namespace Sessions;

/// <summary>
/// A person giving one or more sessions.
/// </summary>
/// <param name="Id">Stable identifier for this speaker.</param>
/// <param name="Name">The name to print in the program.</param>
/// <param name="Bio">A sentence or two, if the speaker sent any.</param>
public sealed record Speaker(int Id, string Name, string? Bio);
\end{minted}

\index{GraphQLIgnore@\code{[GraphQLIgnore]}}%
\code{SpeakerId} is the one line here that is a decision rather than a
transcription. It is how a session finds its speaker in memory, and it has no
business being in the schema: a client that reads it is reading a storage
detail, and a client that reads it will eventually be a client that depends on
it. \code{[property: GraphQLIgnore]} keeps it out.
Section~\ref{sec:ch02-what-the-schema-says-about-your-types} is about why that
attribute has to exist at all.

\code{ConferenceData.cs} holds four sessions and three speakers, chosen so that
one speaker gives two sessions and one speaker sent no biography. Both of those
turn up again before the chapter ends.

\begin{minted}{csharp}
namespace Sessions;

/// <summary>
/// The conference program, held in memory. The next chapter puts it in
/// SQLite with EF Core; nothing outside this file knows where the data
/// lives.
/// </summary>
public static class ConferenceData
{
    public static readonly IReadOnlyList<Speaker> Speakers =
    [
        new Speaker(1, "Ada Fischer", "Works on query planners."),
        new Speaker(2, "Bruno Kaminski", null),
        new Speaker(3, "Chidi Okafor", "Runs platform engineering.")
    ];

    public static readonly IReadOnlyList<Session> Sessions =
    [
        new Session(
            1,
            "Schemas That Outlive Their Authors",
            "A five-year-old schema, and the decisions that shaped it.",
            1,
            new DateTimeOffset(2026, 9, 14, 9, 0, 0, TimeSpan.Zero),
            45),
        new Session(
            2,
            "Reading a Query Plan",
            null,
            1,
            new DateTimeOffset(2026, 9, 14, 10, 0, 0, TimeSpan.Zero),
            45),
        new Session(
            3,
            "Paging Without Offsets",
            "Why the page you get is not the page you asked for.",
            2,
            new DateTimeOffset(2026, 9, 14, 11, 0, 0, TimeSpan.Zero),
            30),
        new Session(
            4,
            "The Bank That Split Its Graph",
            "Eighteen months, four teams, and the unwritten parts.",
            3,
            new DateTimeOffset(2026, 9, 14, 14, 0, 0, TimeSpan.Zero),
            60)
    ];
}
\end{minted}

\code{Query.cs} replaces the hello-world of
section~\ref{sec:ch02-the-smallest-thing-that-runs} entirely.

\begin{minted}{csharp}
namespace Sessions;

[QueryType]
public static class Query
{
    /// <summary>
    /// Every session on the schedule, earliest first.
    /// </summary>
    public static IEnumerable<Session> GetSessions()
        => ConferenceData.Sessions.OrderBy(s => s.StartsAt);

    /// <summary>
    /// One session by its identifier, or null if no session has it.
    /// </summary>
    public static Session? GetSessionById(int id)
        => ConferenceData.Sessions.FirstOrDefault(s => s.Id == id);
}
\end{minted}

\index{resolver!on a type extension}%
That leaves the field that is not a property of anything. A session's speaker
lives in a different collection, so it is computed rather than stored, and it
goes in \code{SessionType.cs}, on a class that extends \code{Session} the same
way \code{Query} extends the root type.

\begin{minted}{csharp}
namespace Sessions;

[ObjectType<Session>]
public static partial class SessionType
{
    /// <summary>
    /// The person giving this session.
    /// </summary>
    public static Speaker? GetSpeaker([Parent] Session session)
        => ConferenceData.Speakers
            .FirstOrDefault(s => s.Id == session.SpeakerId);
}
\end{minted}

\index{Parent@\code{[Parent]}}%
\index{resolver!definition}%
\code{[Parent]} is the whole of the mechanism. It marks the parameter that
receives the object this field is being resolved on, so the method is an
ordinary static function taking a \code{Session} and returning a
\code{Speaker}.
A method that computes one field's value is what GraphQL calls a resolver,
and every method in this chapter that reaches the schema is one;
\code{GetSpeaker} is the first that has to be handed something to work from.
That is worth pausing on, because it is the same shape you will
write in chapter~\ref{ch:first-subgraph} to resolve an entity the router asks
for by key, and the same shape that quietly issues one query per parent in
chapter~\ref{ch:router-n-plus-1}. This one runs against a list in memory, so it
costs nothing here. It will not stay free.

\code{Program.cs} does not change at all. It registers whatever the generator
found, and the generator has now found \code{SessionType} as well as
\code{Query}, so the same call registers more than it did in
section~\ref{sec:ch02-the-smallest-thing-that-runs}.

Run it, and ask for something that spans both records.

\begin{minted}{text}
POST /graphql HTTP/1.1
Host: localhost:5001
Content-Type: application/json

{"query":"{ sessions { id title startsAt speaker { name bio } } }"}
\end{minted}

and back:

\begin{minted}[fontsize=\footnotesize]{json}
{"data":{"sessions":[
  {"id":1,"title":"Schemas That Outlive Their Authors",
   "startsAt":"2026-09-14T09:00:00Z",
   "speaker":{"name":"Ada Fischer","bio":"Works on query planners."}},
  {"id":2,"title":"Reading a Query Plan",
   "startsAt":"2026-09-14T10:00:00Z",
   "speaker":{"name":"Ada Fischer","bio":"Works on query planners."}},
  {"id":3,"title":"Paging Without Offsets",
   "startsAt":"2026-09-14T11:00:00Z",
   "speaker":{"name":"Bruno Kaminski","bio":null}},
  {"id":4,"title":"The Bank That Split Its Graph",
   "startsAt":"2026-09-14T14:00:00Z",
   "speaker":{"name":"Chidi Okafor","bio":"Runs platform engineering."}}]}}
\end{minted}

The response above is the server's, re-indented so that it fits the measure;
the server sends it on one line. Ada Fischer appears twice, once for each of
her sessions, and Bruno Kaminski's \code{bio} came back as \code{null} rather
than as an omitted key. Both of those are the schema doing what the C\# told it
to, and the next section is where you can read what that was.
